\section{What remains to decide}\label{sec:roadmap}

The results leave two questions: which coded updates can recover the required
outputs, and whether those updates can be executed in linear time on tapes.
The local-coding theorems exclude specific information-access patterns. The
motion theorems exclude specific execution patterns. The padding construction
supplies the cut condition, and for linear-time payloads the full limiting
conditions, without requiring a transpose-specific computation.

\paragraph{An impossibility target.}
The earlier paper proves that \textup{(T)} gives causal complete-exchange codes
on valid tape diagrams, at the prescribed block capacities, on input sets with
cardinality deficit at most $Cd$ for a fixed $C$. The diagrams append their
output in the last $P$ steps. It therefore suffices to show that, for every
fixed choice of resource constants and $C$, arbitrarily large dimensions admit
no such diagram and code. This requires less than the full robust Tape Coding Conjecture, which
excludes a larger class of input sets at every sufficiently large size.
\upformal{Transpose.Tape.Statements.Coding}{Transpose.Tape.codes_coentLinear_of_linearTimeTransposable,Transpose.Tape.not_linearTimeTransposable_of_coentLinear_end_io}

Applying our local-coding bounds to this target requires a structural
argument: one must prove that the extracted tape codes satisfy the relevant
locality or factorization hypotheses. The packet theorem already permits a
globally recoded bank. Its bounded-degree consequence applies to fixed
two-coordinate packets after a full-cube bijection, with bounded stored-row
degree, inverse support and reuse, and sublinear combined input deficit and
helper entropy. Fast tape execution alone does not supply these properties.
Likewise, an evolving coded bank is not an original row, and an
arbitrary execution need not factor into the two stages of our helper theorem.
Improving constants in the local bounds does not supply this missing argument.

\paragraph{A lower-bound invariant to test.}
A further lower bound must use more than the aggregate cut and $M^*$
conditions supplied by the padding wrapper. One candidate is an invariant
that tracks both causal availability of coded information and its use in
recovering specified outputs. A useful test is whether adjoining irrelevant
padding leaves the invariant unchanged. The next tests are whether it survives
the correlated and nonlinear examples allowed by our local theorems, and
whether its accumulated cost is superlinear in the full tape model.
This is a proposed direction; we have not proved such an invariant exists.

\paragraph{A construction target.}
The next constructive step is an update and decoding identity with a proved
saving. Its state invariant must specify what is stored, which data each
update can access, and how the next outputs are recovered. The cost must
include helper and layout generation, encoding, decoding, head movement,
finite control and raw output. A saving in an abstract coding model alone
does not establish \textup{(T)}.

The first paper's code-gap theorem requires an $\Omega(\log d)$ advantage over
routing on the relevant tape networks. A code with this gain would still need
a uniform finite-state implementation with linear total cost.
\upformal{Transpose.Tape.Statements.Coding}{Transpose.Tape.code_gap}

Several possibilities remain beyond the stated corollaries: high-treewidth queries
on original rows that do not fit the two-coordinate packet model; larger
packet supports; global recodings with unbounded inverse support or reuse;
more than two nonlinear stages; and recursive updates that fail the stated
affine-regularity hypotheses. A preliminary global bijection with bounded
inverse support and reuse is already covered: on bounded-degree packet
graphs it cannot produce a fixed fractional broadcast saving with
$D_w+H(Z)=o(P)$. None of the remaining possibilities is thereby shown feasible.
Larger inverse supports must still be tested against the finite packet bound.
The padding construction supplies motion, not an encoder or decoder for
complete exchange.

\paragraph{Evaluating a candidate.}
State its information access, helper costs, correctness set and actual tape
costs first. Then evaluate the applicable quantitative bounds. Discard the
candidate when a bound contradicts its claimed cost; membership in a local
model alone is not enough. If no exclusion applies, identify the exact
hypothesis it escapes and test the remaining computation directly.
